% Root and agent B passed the complete literal proof; the control certificate has an independent gate.
\subsection{Connected five-state graphs opposite a right-zero corner}

\begin{lemma}[The natural connected-graph Tor comparison]
\label{C:connected-graph-cone-natural-Tor}
Let a finite monoid $M$ act cellularly on the right on a connected
finite graph $\Gamma$. Put
\[
 X=\operatorname{hocofib}(\Gamma_{hM}\longrightarrow BM),
 \qquad W=H_1(\Gamma;\mathbb Z).
\]
Then $H_1(X;\mathbb Z)=0$ and, naturally in the graph and action,
\[
 \widetilde H_n(X;\mathbb Z)
   \cong\operatorname{Tor}_{n-2}^{\mathbb ZM}(W,\mathbb Z),
       \qquad n\geq2.
\]
The analogous formula holds over every field.
If the action factors through a surjective monoid map $M\to M'$, the
actual cofiber comparison is an isomorphism on $H_2$ and a surjection
on $H_3$, integrally and over every field.
If an $M$-equivariant self-map of $\Gamma$ acts on $W$ as
$d\,1_W$ for an integer $d\notin\{0,1\}$, every positive rational
cohomology product of $X$ vanishes.
\end{lemma}
\begin{proof}
The augmented cellular complex
$C_1(\Gamma)\to C_0(\Gamma)\to\mathbb Z$ is a complex of
right $\mathbb ZM$-modules. Its only nonzero homology is $W$
in degree one, and the literal inclusion of cycles is an
equivariant quasi-isomorphism from $W$ in that degree.
Tensor with a free left bar resolution of the trivial module.
The mapping cone of the augmentation computes the homotopy
cofiber and supplies the further degree shift. This gives the
displayed formula and its naturality, including the field version.

For a quotient of the acting monoid, the degree-zero two-sided bar
is the identical module $W$, and its degree-one map is surjective.
The degree-zero differentials agree, so every target one-cycle lifts
to a source one-cycle. This proves the $H_2$ isomorphism and $H_3$
surjection in the stated directions.
For a scalar self-map, the literal map on graph cycles is $d$;
the same natural chain comparison makes its induced map on every
positive-degree rational cohomology group multiplication by $d$.
Multiplicativity gives $d(xy)=d^2xy$, hence $xy=0$ over $\mathbb Q$.
\end{proof}

\begin{theorem}[The connected graph case at five cross states]
\label{C:right-zero-five-connected}
Let $R=\{1,p,q\}$ with $ab=b$ for $a,b\in\{p,q\}$, and let $S$
be an $(R,M)$-biset with at most five elements. Write $P,Q:S\to U$
for the two retractions onto their common image, and form the
natural action graph $\Gamma$ of
Lemma~\ref{C:right-zero-action-graph}: its vertices are $U$, and
each $s\in S\setminus U$ is an oriented edge from $Ps$ to $Qs$.
If this graph is connected, the two-object collage cannot have
classifying space $\mathbb CP^2$. There is no bound on $|M|$.
\end{theorem}
\begin{proof}
The graph comparison is natural for the right $M$-action, and
$BR$ is contractible. Thus the collage is the homotopy cofiber
in the preceding lemma. A vertex fixed by all of $M$ gives zero
degree-two products by the opposite form of
Lemma~\ref{C:fixed-incoming-cups}, since $H^2(BR;\mathbb Q)=0$.
An actor inducing zero on $W$ makes its coinvariants zero and
therefore makes $H_2(X;\mathbb Q)=0$.
We may exclude these possibilities throughout.

A connected graph on at most five total vertices and edges has
at most three vertices. With one vertex it has the fixed vertex
just discussed; with three vertices it is a tree, so its map to
a point is an equivariant underlying equivalence and the cofiber
is contractible. With two vertices and at most two edges, the
connected cases are precisely those treated in the proof of
Theorem~\ref{C:right-zero-four-state-reduction}: a tree, a loop
and a connecting edge, or two parallel connecting edges. Their
fixed-vertex, zero-degree and circle-power arguments exclude a
projective plane without any bound on the source monoid.

It remains that $|U|=2$ and there are three edges. Label the
vertices $0,1$ and let $k$ be the number of edges joining distinct
vertices. Since the graph is connected, $k\in\{1,2,3\}$.
Every actor maps an edge to an edge of the corresponding ordered
endpoint type or collapses it to a vertex. In particular it cannot
reverse an edge while fixing its endpoints. We enumerate the
following possibilities up to reversing all orientations and
interchanging the vertices.

\emph{Three connecting edges.}
If all have the same orientation, no actor can interchange the
vertices. A vertex-constant actor kills all graph homology;
otherwise both vertices are fixed. Thus the only remaining shape
has two edges $a,b:0\to1$ and one edge $c:1\to0$.
A vertex-constant actor again kills all graph homology, while in
the absence of a vertex interchange both vertices are fixed.
Hence every actor fixes or interchanges the vertices, with at least
one interchange. The quotient identifying $a$ with $b$ is strictly
$M$-equivariant and is the two-edge circle. Write
$z_a=a+c$, $z_b=b+c$ for the integral cycle basis.
An interchanging actor sends both basis cycles to the same one,
so their difference is zero in coinvariants. Every actor acts
as the identity on the quotient cycle. The induced map on cycle
coinvariants is therefore an isomorphism to $\mathbb Z$.
The natural Tor lemma makes the cofiber comparison an isomorphism
on second rational cohomology in the pullback direction.
On the quotient circle every actor is the identity or the
half-rotation. The degree-three map commutes with both, so its
cofiber has zero rational products by the scalar-map clause.
Pullback excludes the original nonzero square.

\emph{Two connecting edges and one loop.}
Put the loop at vertex zero. If the connecting edges have the
same orientation, a vertex interchange is impossible. A constant
map to vertex one kills all graph homology; in its absence vertex
zero is fixed. Thus the connecting edges must have opposite
orientations. If there is no interchanging actor, the same
alternative applies. An interchanging actor collapses the loop,
so its class is zero in cycle coinvariants. A vertex-constant
actor sends the connecting cycle to a multiple of this loop
class or to zero; it would then kill the remaining coinvariant
class as well. Consequently all actors fix or interchange the
vertices. Collapsing the loop is now a strictly equivariant map
to the two-edge circle. It induces an isomorphism on the
coinvariants of graph homology, since the loop class was already
killed and the remaining circle class has trivial action.
Use the same natural comparison and degree-three circle map.

\emph{One connecting edge and two loops.}
No actor can interchange the vertices, because no reverse
connecting edge exists. If both loops are based at one vertex,
an actor constant at the other vertex kills all graph homology;
in its absence the loop vertex is fixed. It remains that there
is one loop $\ell_i$ at each vertex $i$, and a connecting edge
$a:0\to1$.

To avoid a common fixed vertex, the action image must contain
actors constant at each of the two vertices. A nonzero invariant
functional on the cycle coinvariants forces both loop classes
to have the same nonzero value. Indeed a constant-at-$i$ actor
sends both loops either to $\ell_i$ or to its vertex; if any
such loop image is collapsed, the relations from constant actors
at both vertices force both functional values to vanish.
Likewise a vertex-fixing actor must fix both loops. Thus every
vertex-fixing actor acts as the identity on the entire graph,
and every other actor has exactly the form
\[
 c_{i,\delta}:\quad
 0,1\longmapsto i,\qquad
 \ell_0,\ell_1\longmapsto\ell_i,\qquad
 a\longmapsto
 \begin{cases}i,&\delta=0,\\\ell_i,&\delta=1.
 \end{cases}
\]
Here $i,\delta\in\{0,1\}$.

If all such actors have $\delta=0$, fix the connecting edge and
map each loop to its square. This is a strictly equivariant
graph self-map acting as twice the identity on graph homology.
If all have $\delta=1$, map each loop to its square and map the
connecting edge to the path $\ell_0a$, using two equal parameter
subintervals. A constant-at-$i$ actor sends this path to
$\ell_i^2$, exactly the image required by equivariance; the
vertex-fixing actors are identities. This again gives a scalar
degree-two self-map. Both uniform cases have zero rational products.

If both values of $\delta$ occur, closure and the presence of both
vertex values force all four maps $c_{i,\delta}$.
For the right-action convention their multiplication is
\[
 c_{i,\delta}c_{j,\epsilon}=c_{j,\delta}.
\]
The faithful image of $M$ is therefore exactly the identity
adjoined to this four-element rectangular band. The original
collage maps by an actual functor to the thirteen-arrow collage
with this five-element source image, the same corner $R$, and
the same five states. The graph-Tor comparison makes this
map an isomorphism on $H_2(-;\mathbb F_2)$.

The literal target table and a normalized two-cocycle and square
primitive are stored in
\path{code/outputs/a_rightzero_five_connected_control.json}.
Its normalized chain dimensions are $11,50,212,872$ in degrees
one through four. Exact elimination gives $H_1(-;\mathbb F_2)=0$
and $\dim H^2(-;\mathbb F_2)=1$; the supplied generator is not
a coboundary, and all $872$ equations verify its square primitive.
The standalone checker \path{code/b_verify_rightzero_five_control.py}
reconstructs the two retractions, all right maps, and the entire
category table without producer imports. It checks every typed
associativity equation, all consecutive-boundary equations, the
nonboundary generator and all literal square identities. Its complete
record is
\path{code/outputs/b_rightzero_five_connected_control_verified.json}.
Thus its degree-two square is zero. Naturality and the actual
second-homology isomorphism make the square zero on the original
collage, contradicting the mod-two ring of $\mathbb CP^2$.
These graph types exhaust the connected case.
\end{proof}

\begin{remark}[Remaining five-state graphs]
\label{C:right-zero-five-connected-scope}
This theorem concerns the connected graph of the left $R$-action,
not merely a connected action category after taking the right
$M$-bar. A disconnected graph can have its components joined by
the right action. Those cases are handled separately by
Theorem~\ref{C:right-zero-five-state-reduction} below.
The thirteen-arrow certificate treats one faithful action image;
the natural graph-Tor comparison is what extends that calculation
to arbitrary finite source monoids with that image.
\end{remark}
